\documentclass[10pt,a3paper,landscape]{article}
\usepackage[margin=6mm]{geometry}
\usepackage[T1]{fontenc}
\usepackage[utf8]{inputenc}
\usepackage[french]{babel}
\usepackage{lmodern}
\usepackage{microtype}
\makeatletter\def\input@path{{../../../Style/}}\makeatother
\NeedsTeXFormat{LaTeX2e}
\ProvidesPackage{TLPosterCarteMentale}[2026/08/03 Posters et cartes mentales SI]

\RequirePackage{tikz}
\RequirePackage[most]{tcolorbox}
\RequirePackage{xcolor}
\RequirePackage{amsmath,amssymb}
\RequirePackage{graphicx}
\usetikzlibrary{positioning,arrows.meta,calc,shapes.geometric,mindmap,backgrounds,fit}

\definecolor{TLPosterBlue}{RGB}{24,86,141}
\definecolor{TLPosterLight}{RGB}{234,243,250}
\definecolor{TLPosterAccent}{RGB}{232,126,34}
\definecolor{TLPosterGreen}{RGB}{52,152,102}
\definecolor{TLPosterRed}{RGB}{192,57,43}
\definecolor{TLPosterGray}{RGB}{80,90,100}

\newcommand{\TLPosterTitre}[2]{%
  \begin{tcolorbox}[enhanced,boxrule=0pt,arc=0pt,colback=TLPosterBlue,left=6mm,right=6mm,top=4mm,bottom=4mm]
    {\color{white}\bfseries\LARGE #1}\hfill{\color{white}\large #2}
  \end{tcolorbox}%
}

\newtcolorbox{TLPosterBloc}[2][]{
  enhanced,breakable,
  colback=TLPosterLight,
  colframe=TLPosterBlue,
  colbacktitle=TLPosterBlue,
  coltitle=white,
  fonttitle=\bfseries,
  title={#2},
  boxrule=0.8pt,
  arc=2mm,
  left=3mm,right=3mm,top=2mm,bottom=2mm,
  #1
}

\newtcolorbox{TLPosterFormule}[1][]{
  enhanced,
  colback=white,
  colframe=TLPosterAccent,
  boxrule=1pt,
  arc=2mm,
  fontupper=\bfseries,
  halign=center,
  #1
}

\newcommand{\TLPosterMethode}[1]{%
  \begin{tcolorbox}[enhanced,colback=TLPosterGreen!8,colframe=TLPosterGreen,title=Méthode,fonttitle=\bfseries]
  #1
  \end{tcolorbox}%
}

\newcommand{\TLPosterAttention}[1]{%
  \begin{tcolorbox}[enhanced,colback=TLPosterRed!6,colframe=TLPosterRed,title=Point de vigilance,fonttitle=\bfseries]
  #1
  \end{tcolorbox}%
}

\tikzset{
  TLmindroot/.style={concept,concept color=TLPosterBlue,text=white,font=\bfseries\large,minimum size=34mm},
  TLmindlevel1/.style={level 1 concept/.append style={font=\bfseries,minimum size=23mm,sibling angle=45}},
  TLmindlevel2/.style={level 2 concept/.append style={font=\small,minimum size=17mm}},
  TLarrow/.style={-{Stealth[length=3mm]},thick,draw=TLPosterBlue},
  TLbox/.style={rounded corners=2mm,draw=TLPosterBlue,fill=TLPosterLight,align=center,inner sep=3mm}
}

\newenvironment{TLCarteMentale}[1]{%
  \begin{tikzpicture}[mindmap,grow cyclic,concept color=TLPosterBlue,every node/.style={concept},TLmindlevel1,TLmindlevel2]
  \node[TLmindroot]{#1}
}{;
  \end{tikzpicture}
}

\newcommand{\TLBranche}[3][]{child[concept color=#2]{node[#1]{#3}}}

\endinput

\pagestyle{empty}
\renewcommand{\familydefault}{\sfdefault}
\begin{document}
\begin{tikzpicture}
\TLMapBase{Résistance des matériaux}{Charges $\rightarrow$ Cohésion $\rightarrow$ Sollicitation $\rightarrow$ Contraintes / déformée $\rightarrow$ Vérifier}
\TLMapBranchTL{À connaître}{\textbullet\ poutre élancée, petites déformations, élasticité linéaire\\[1mm]\textbullet\ Saint-Venant et Navier--Bernoulli\\[1mm]\textbullet\ cohésion : $N,T_y,T_z,M_t,M_{fy},M_{fz}$\\[1mm]\textbullet\ traction : $\sigma=N/S$, $\sigma=E\varepsilon$\\[1mm]\textbullet\ flexion : $\sigma=-M_fy/I$, $EIy''=M_f$\\[1mm]\textbullet\ résistance ET rigidité}
\TLMapBranchTR{Méthode cohésion}{1. Calculer les réactions d'appui.\\[1mm]2. Définir les intervalles.\\[1mm]3. Couper en $x$.\\[1mm]4. Isoler le côté avec le moins de charges.\\[1mm]5. Réduire les actions en $G$.\\[1mm]6. Appliquer le PFS et identifier $N,T,M_t,M_f$.\\[1mm]7. Tracer les diagrammes ; contrôler sauts et continuités.}
\TLMapBranchML{Cas traction / compression}{\textbf{Reconnaître :} torseur réduit à $N$.\\[1mm]\textbf{Résistance :} $\sigma=N/S$ ; avec singularité $\sigma_{max}=K_t|N|/S_{nette}$ ; comparer à $R_e/s$.\\[1mm]\textbf{Déformation :} $\varepsilon=\sigma/E$ ; si constantes, $\Delta L=NL/(ES)$.\\[1mm]\textbf{Signe :} $N>0$ traction, $N<0$ compression.}
\TLMapBranchMR{Cas flexion}{1. Déterminer $T(x)$ et $M_f(x)$.\\[1mm]2. Localiser $|M_f|_{max}$.\\[1mm]3. Résistance : $\sigma=-M_fy/I$ ou $|\sigma_{max}|=|M_f|/W$.\\[1mm]4. Rigidité : intégrer $EIy''=M_f(x)$.\\[1mm]5. Appliquer conditions aux limites / continuité.\\[1mm]6. Comparer $f_{max}$ à $f_{limite}$.}
\TLMapBranchBL{Hypothèses / pièges}{\textbullet\ vérifier $L\gg$ dimension de section\\[1mm]\textbullet\ ne pas appliquer Saint-Venant près d'une charge locale ou d'un accident géométrique\\[1mm]\textbullet\ la superposition exige un comportement linéaire élastique\\[1mm]\textbullet\ une pièce peut résister mais être trop souple\\[1mm]\textbullet\ ne pas oublier les concentrations de contrainte}
\TLMapBranchBR{Décision finale}{\textbf{Si $N$ :} traction/compression.\\[1mm]\textbf{Si $M_f$ :} flexion.\\[1mm]\textbf{Si plusieurs composantes :} partir du torseur, traiter les effets compatibles avec le modèle et superposer seulement si autorisé.\\[1mm]\textbf{Toujours :} sécurité $s$, résistance, rigidité, unités, ordre de grandeur.}
\TLMapRibbon{Une formule de contrainte n'est valable qu'après identification de la sollicitation par le torseur de cohésion.}
\end{tikzpicture}
\end{document}
